\section[Notebook 1972-9: a finite extension]{Notebook 1972-9: $k[t, x]$ is finite over $k[t, tx + x^q]$\protect\verifie{Notebook1972_9Finite}}\label{sec:quillen-1972-9-finite}

\emph{Context.} On page~71 of \carnet{1972}{9} Quillen deduces that the family $f + t\cdot\mathrm{id}$ is proper from the finiteness of $k[t, x]$ over the subalgebra generated by $t$ and $tx + x^q$.

\begin{encadre}
\begin{theoreme}[Quillen, notebook 1972-9, p.~71]\label{thm:quillen-1972-9-finite}
Let $k$ be a commutative ring and $q \ge 2$. Then $k[t, x]$ is a finite module over its subalgebra $k[t, y]$, $y = tx + x^q$.
\end{theoreme}
\end{encadre}

\begin{proof}
$x$ is a root of the monic polynomial $Z^q + tZ - y$ with coefficients in $k[t, y]$, so $k[t, y][x]$ is finite over $k[t, y]$; and $k[t, y][x] = k[t, x]$, since $t$ lies in $k[t, y]$.
\end{proof}

\emph{Verification in Lean.} The file \href{\leanblobmain/Notebook1972_9Finite.lean}{\texttt{Notebook1972\_9Finite.\allowbreak{}lean}} proves \texttt{isIntegral\_X\_one} and \texttt{finite}, in \texttt{MvPolynomial (Fin 2) k}. The hypothesis $q \ge 2$ makes the polynomial monic; for $q = 1$ the statement fails ($y = (1 + t)x$), which is not formalised. The axioms used are the three usual axioms of mathlib.

\litterature{Elementary (integral extensions); the edition does not present it as new.}
